\documentstyle[12pt]{article}
\setlength{\oddsidemargin}{0in}
\setlength{\evensidemargin}{-0in}
\setlength{\textwidth}{6.4 in}
\setlength{\topmargin}{-.8in}
\setlength{\textheight}{9.2in}

\date{}
\title{{\bf Edward Chang\\Statement of Teaching Interests}}
\author{}
\begin{document}
\maketitle
\vspace{-2.0cm}

\def\thepage{}

\vspace{0.3cm}
%The stories of great teachers like Plato and Confucius
%have inspired me to teach since my formative years.
I believe that being a good teacher requires not only the desire to
pass knowledge on to others, but also the ability to
cultivate and inspire young minds to discover the unknown 
and to expand human possibilities. 
Teaching is a privilege that also carries 
great respoisibilities. At the very least, a teacher must 
present the course materials clearly.  More important,
a teacher must motivate the interest
of students with different backgrounds and levels of
maturity and must inspire them to progress
beyond what they think they can do.

\vspace{0.3cm}
To prepare myself to teach effectively, I have actively sought out
opportunities to improve my teaching skills.
As a graduate student, I was a teaching assistant for
one undergraduate course (Introduction to Programming Languages)
and one graduate course (Advanced Topics in Operating Systems).
I prepared and graded assignments and solutions, and assisted 
in preparing and grading exams. I also led discussion 
sessions, which proved a useful way to learn to address the 
students' weak points directly. 

\vspace{0.3cm}
In addition, I had significant teaching experience as
the instructor of a database system course in the
summer of 1998 at Stanford. I conducted six full lectures,
and I also led 20-minute discussion sessions after each of $12$
taped lectures given by professor Ullman in an earlier quarter. 
The course was broadcast over Stanford's Instructional
Television Network. Interacting with forty students
both local and remote,
and with both full-time and part-time students working in Silicon
Valley companies presented challenging problems and 
opportunities because the students' backgrounds, motivation, 
and dedication varied greatly. 
Directing this class was a wonderful experience. In particular,
I enjoyed the interaction with students, which gave me
the chance to learn about their needs and to assist them on an
individual basis. In this course, I also used the first draft of a new 
textbook written by professors Garcia-Molina, Ullman, 
and Wildom; I gathered and provided feedback to the 
professors for revisions before the first edition came out.  
I received very 
positive evaluations from both the professors and the students.
My average rating in the course evaluation was $4.1$ out of $5$, 
where $4$ is ``very good'' and $5$ is ``excellent.''
(The students' evaluations are available upon request.)

\vspace{0.3cm}
I believe that my experience has prepared me for the teaching
requirements of an academic position. Moreover, 
my desire to teach well and my patience should
help me remain a very good teacher. 
I am interested in teaching both
undergraduate and graduate level classes in databases, operating
systems, distributed systems, computer architectures, multimedia
systems, programming
languages, and other related courses in Computer Science
and Electrical Engineering disciplines. I look forward to 
having a positive impact on students.

\end{document}
